Avendo individuato i costrutti rispetto ai quali l’insieme di tutti gli insiemi regolari resta invariato (concatenazione, unione e operatore * di Kleene), si può formalizzare una grammatica data dalla rappresentazione
simbolica di alcuni insiemi regolari di base.
\begin{defi}
    Fissiamo un alfabeto $A$. L'insieme $\mathcal E$ (anche indicato con $E$) delle \textit{espressioni regolari} sull'alfabeto 
$A$ è il più piccolo insieme tale che: 
\begin{itemize}
    \item $\mathbb  0, \mathbb 1 \in \mathcal E$;

    \item  per ogni $a \in A$ $a \in \mathcal E$;

    \item se $E_1, E_2 \in \mathcal E $ allora 
    \(
    E_1 + E_2,
    E_1 \cdot E_2,
    E_1^*\in \mathcal E
    \)
\end{itemize}
\end{defi}
\begin{exam}
    Alcune espressioni regolari sull’alfabeto $A= {0,1}$ sono:
    \begin{itemize}
        \item $\mathbb 0$;

        \item $\mathbb 0+0$;

        \item $(0+1)^*$;

        \item $0\cdot (0+1\cdot 1)^*\cdot \mathbb 1$.
    \end{itemize}
\end{exam}
\begin{rema}
    Chiaramente, fino ad ora, abbiamo usato i simboli $+,\cdot, \cdot$ in due contesti diversi: per la definizione formale di espressione regolare e per le definizioni di operazioni sui sottoinsiemi degli alfabeti.
    Questo non è un problema visto che si capisce a quale definizione si fa riferimento in base al contesto.
\end{rema}
\begin{defi}[Interpretazione di un'espressione regolare]
\label{def:22}
Per ogni espressione regolare $e\in E$, di alfabeto $A$, definiamo un'interpretazione
\[
\llbracket\ \rrbracket : E \longrightarrow 2^{A^*}
\]
dove
\[
\llbracket e\rrbracket :=
\left\{
\begin{matrix}
\varnothing & \text{se } e=\mathbb{0} & \text{(insieme vuoto)}\\[2pt]
\{\epsilon\} & \text{se } e=\mathbb{1} & \text{(insieme singleton parola vuota)}\\[2pt]
\{a\} & \text{se } e=a\in A & \text{(insieme singleton simbolo)}\\[2pt]
\llbracket e_1\rrbracket\cdot\llbracket e_2\rrbracket & \text{se } e=e_1\cdot e_2 & \text{(concatenazione)}\\[2pt]
\llbracket e_1\rrbracket+\llbracket e_2\rrbracket & \text{se } e=e_1+e_2 & \text{(unione)}\\[2pt]
\llbracket e_1\rrbracket^* & \text{se } e=e_1^* & \text{(insieme star)}
\end{matrix}
\right.
\]
\end{defi}
\begin{figure}[H]
\centering
\begin{tikzpicture}[
    >=stealth', shorten >=1pt, auto,
    every state/.style={circle, draw, fill=white, minimum size=0.8cm, inner sep=1pt}
]
% Rettangolo di A
\draw[-, line width=2.5pt, rounded corners=3pt] (2,-1.4) rectangle (7,1.4);
\node at (4.5,0) {$\mathcal{A}$};

% Stati (q_0 ora sul bordo sinistro, in sovrimpressione)
\node[state, accepting] (q0p) at (0,0)    {$q_0'$};
\node[state]            (q0)  at (2,0)    {$q_0$};
\node[state, accepting] (qi)  at (7,0.8)  {$q_i$};
\node[state, accepting] (qj)  at (7,-0.8) {$q_j$};

% Freccia d'ingresso e epsilon-mossa iniziale
\draw[->] ([shift={(-0.7,0.7)}]q0p.center) -- (q0p);
\draw[->] (q0p) -- node {$\varepsilon$} (q0);

% epsilon-mossa di ritorno interna (da q_i)
\draw[->, rounded corners=8pt]
    (qi.45) -- (qi.45 |- 0,1.9) -- (q0.100 |- 0,1.9) -- (q0.100);
\node at (7.55,1.5) {$\varepsilon$};

% epsilon-mossa di ritorno esterna (da q_j)
\draw[->, rounded corners=8pt]
    (qj.0) -- (8.1,-0.8) -- (8.1,2.4) -- (q0.135 |- 0,2.4) -- (q0.135);
\node at (8.4,-0.4) {$\varepsilon$};
\end{tikzpicture}
\caption{Costruzione dell'automa che decide $X^*$ a partire da quello che decide $X$.}
\label{fig:kleene}
\end{figure}

\begin{rema}
    Chiaramente, riprendendo la precedente osservazione, con $\llbracket e_1\rrbracket\cdot\llbracket e_2\rrbracket$ indichiamo la conatenazione degli insiemi $\llbracket e_1\rrbracket$ e $\llbracket e_2\rrbracket$.
\end{rema}

\begin{prop} Data un’espressione regolare $e\in E$ l’insieme $\llbracket e \rrbracket$ è un insieme regolare.
\begin{proof} Questo risultato è conseguenza delle proprietà di chiusura (per prodotto, somma e star)
soddisfatte dagli insiemi DFA-decidibili, e dal fatto che le tre espressioni regolari di base si interpretano
con insiemi DFA-decidibili: si veda la figura 
\ref{fig:automi-base} per gli automi che decidono le tre espressioni regolari di base.
\end{proof}
\end{prop}
\begin{figure}[H]
\centering
\begin{subfigure}[t]{0.3\textwidth}
    \centering
    \begin{tikzpicture}[
        ->, >=stealth', shorten >=1pt, auto, node distance=3cm,
        every state/.style={circle, draw, thick, minimum size=0.8cm, inner sep=1pt}
    ]
    \node[state] (q0) {$q_0$};
    \draw[->] ([shift={(-0.8,0.8)}]q0.center) -- (q0);
    \end{tikzpicture}
    \caption{NDFA che decide l'insieme vuoto}
    \label{fig:ndfa-vuoto}
\end{subfigure}
\hfill
\begin{subfigure}[t]{0.3\textwidth}
    \centering
    \begin{tikzpicture}[
        ->, >=stealth', shorten >=1pt, auto, node distance=3cm,
        every state/.style={circle, draw, thick, minimum size=0.8cm, inner sep=1pt}
    ]
    \node[state, accepting] (q0) {$q_0$};
    \draw[->] ([shift={(-0.8,0.8)}]q0.center) -- (q0);
    \end{tikzpicture}
    \caption{NDFA che decide l'insieme singleton della parola vuota}
    \label{fig:ndfa-epsilon}
\end{subfigure}
\hfill
\begin{subfigure}[t]{0.36\textwidth}
    \centering
    \begin{tikzpicture}[
        ->, >=stealth', shorten >=1pt, auto, node distance=2.2cm,
        every state/.style={circle, draw, thick, minimum size=0.8cm, inner sep=1pt}
    ]
    \node[state] (q0) {$q_0$};
    \node[state, accepting] (q1) [right of=q0] {$q_1$};
    \draw[->] ([shift={(-0.8,0.8)}]q0.center) -- (q0);
    \path (q0) edge node {$a$} (q1);
    \end{tikzpicture}
    \caption{NDFA che decide l'insieme singleton di una lettera dell'alfabeto}
    \label{fig:ndfa-lettera}
\end{subfigure}
\caption{Gli automi che decidono le tre espressioni regolari di base}
\label{fig:automi-base}
\end{figure}
\begin{figure}[H]
\centering
\begin{tikzpicture}[
    >=stealth', shorten >=1pt, auto,
    every state/.style={circle, draw, fill=white, minimum size=0.8cm, inner sep=1pt}
]
% Automa A1 (sopra)
\draw[-, line width=2.5pt, rounded corners=3pt] (3,0.2) rectangle (6.6,2.8);
\node at (4.8,1.5) {$\mathcal{A}_1$};
\node[state] (q0a) at (3,1.5) {$q_0$};
\node[state, accepting] (qia) at (6.6,2.3) {$q_i$};
\node[state, accepting] (qja) at (6.6,0.7) {$q_j$};

% Automa A2 (sotto)
\draw[-, line width=2.5pt, rounded corners=3pt] (3,-2.8) rectangle (6.6,-0.2);
\node at (4.8,-1.5) {$\mathcal{A}_2$};
\node[state] (q0b) at (3,-1.5) {$q_0$};
\node[state, accepting] (qib) at (6.6,-0.7) {$q_i$};
\node[state, accepting] (qjb) at (6.6,-2.3) {$q_j$};

% Nuovo stato iniziale
\node[state] (q0p) at (1,0) {$q_0'$};
\draw[->] ([shift={(-0.7,0.5)}]q0p.center) -- (q0p);

% epsilon-mosse
\draw[->] (q0p) -- node[above left, pos=0.5] {$\varepsilon$} (q0a);
\draw[->] (q0p) -- node[below left, pos=0.5] {$\varepsilon$} (q0b);
\end{tikzpicture}
\caption{Costruzione dell'automa che decide $X_1 + X_2$ a partire da quelli che decidono $X_1$ e $X_2$.}
\label{fig:somma}
\end{figure}
Quello che invece vorremmo mostrare ora è che vale anche l'implicazione inversa della proposizione precedente. Per farlo introdurremo una nuova versione di automa.
\begin{defi}
Un \textit{automa a stati finiti non-deterministico generalizzato} di alfabeto $A$ e insieme degli stati $Q$ è una tripla:
\[
(T, q_0, F)
\]
dove
\[
T : Q \times E \to 2^Q,
\]
$q_0 \in Q$ è detto stato iniziale, $F \subset Q$ è detto insieme degli stati finali accettanti.
\end{defi}
L'idea alla base di questo nuovo tipo di automa è
che si può etichettare le transizioni dell'automa
non solo con le lettere dell'alfabeto ma con
qualsiasi espressione costruita a partire dall'alfabeto.
A questo punto, dovremmo però dare di nuovo le definizioni di
algoritmo associato ad un automa generalizzato e di decidibilità.
Usando però la proposizione sopra, possiamo fare ciò in modo furbo:
\begin{rema}
    \label{aut generalizzato a eps. NDFA}
    Dato un automa generalizzato $A=(T,q_0,F)$ con insieme di stati $Q$,
    gli si può associare un $\epsilon$-NDFA nel seguente modo.
    Per ogni coppia $(q,e)$ tale che
    $T(q,e)=p$,
    il risultato precedente fornisce un DFA $A_e=(T^e, q_0^e, F^e)$
    che decide l'iniseme regolare $\llbracket e \rrbracket$.
    A questo punto, possiamo modificare l'automa $A$ in un $\epsilon$-NDFA della forma $A'=(T', q_0', F')$, 
    eliminando tutte le definizioni del tipo
    $T(q,e)=p$ e aggiungendo per ogni coppia $(q,e)$ come sopra:
    \begin{enumerate}
        \item una $\epsilon$-transizione
        che collega $q$ a $q_0^e$, ovvero definendo $T'(q,\epsilon):=q_0$;

        \item le transizioni interne di ogni $A_e$;
        
        \item una $\epsilon$-transizione
        che collega ogni stato accettante di $A_e$ $q_i$ a $p$, ovvero definendo $T'(q_i,\epsilon):=p$ per ogni $q_i\in F^e$.
    \end{enumerate}
\end{rema}
\begin{figure}[H]
\centering
\begin{tikzpicture}[
    >=stealth', shorten >=1pt, auto,
    every state/.style={circle, draw, fill=white, minimum size=0.8cm, inner sep=1pt}
]
%%%%%%%%%%%%%%%%%%%% Automa di partenza %%%%%%%%%%%%%%%%%%%%
% Stati
\node[state] (q) at (3.6,0) {$q$};
\node[state] (p) at (6.6,0) {$p$};

% Transizione T(q,e) = p
\draw[->] (q) -- node {$e$} (p);

%%%%%%%%%%%%%%%%%%%% Freccia di trasformazione %%%%%%%%%%%%%%%%%%%%
\draw[->, line width=1pt] (5.1,-1) -- node[right] {\textit{sostituzione}} (5.1,-2.6);

%%%%%%%%%%%%%%%%%%%% Automa trasformato %%%%%%%%%%%%%%%%%%%%
% Rettangolo di A_e
\draw[-, line width=2.5pt, rounded corners=3pt] (2.4,-6.4) rectangle (7.4,-3.6);
\node at (4.9,-5) {$\mathcal{A}_e$};

% Stati
\node[state] (qb)  at (0,-5)     {$q$};
\node[state] (q0e) at (2.4,-5)   {$q_0^e$};
\node[state] (qi)  at (7.4,-4.2) {$q_i$};
\node[state] (qj)  at (7.4,-5.8) {$q_j$};
\node[state] (pb)  at (10.2,-5)  {$p$};

% epsilon-mossa di ingresso in A_e
\draw[->] (qb) -- node {$\varepsilon$} (q0e);

% epsilon-mosse di uscita dagli stati accettanti di A_e
\draw[->] (qi) -- node[above, sloped] {$\varepsilon$} (pb);
\draw[->] (qj) -- node[below, sloped] {$\varepsilon$} (pb);
\end{tikzpicture}
\caption{Un passo di sostituzione: la transizione $T(q,e)=p$ viene rimpiazzata dal DFA $\mathcal{A}_e$ che decide $\llbracket e \rrbracket$, collegato a $q$ e a $p$ tramite $\varepsilon$-transizioni.}
\label{fig:sostituzione}
\end{figure}
\begin{defi}
    Dato un automa generalizzato $A$, \textit{l'insieme $X \subset A^*$
    deciso dall'automa generalizzato $A$} è l'insieme deciso dall'$\epsilon$-NDFA
    che gli si può associare usando la procedura esposta nell'osservazione \ref{aut generalizzato a eps. NDFA}.
\end{defi}
Ora elenchiamo tre trasformazioni 
che si possono effettuare su un automa generalizzato
che non modicano l'insieme deciso dall'automa stesso.
\begin{itemize}
    \item Setup restituisce l'automa con l'aggiunta di uno stato iniziale $q_0^*$ e di un unico stato finale $q_F^*$ che non hanno archi entranti e su cui non saranno fatte agire le trasformazioni seguenti (Merge e Erase). Il nuovo automa avrà le stesse transizioni dell'automa dato con l'aggiunta di una $\epsilon$-transizione da $q_0^*$ a $q_0$ e di $\epsilon$-transizioni da ogni \textcolor{red}{Queste epsilon mosse non dovrebbero essere delle $\mathbb 1$-mosse visto che stiamo eticchettando con le espressioni.} stato finale $q \in F$ verso $q_F^*$:
    \[
    T'(q, e) := 
    \begin{cases} 
    q_0 & \text{se } q = q_0^* \text{ e } e = \epsilon \\ 
    q_F^* & \text{se } q \in F \text{ e } e = \epsilon \\ 
    T(q, e) & \text{altrimenti.} 
    \end{cases}
    \]

    \begin{figure}[H]
\centering
\begin{tikzpicture}[
    >=stealth', shorten >=1pt, auto,
    every state/.style={circle, draw, fill=white, minimum size=0.8cm, inner sep=1pt}
]
%%%%%%%%%%%%%%%%%%%% Automa di partenza %%%%%%%%%%%%%%%%%%%%
% Rettangolo di A
\draw[-, line width=2.5pt, rounded corners=3pt] (0,-1.4) rectangle (3,1.4);
\node at (1.5,0) {$\mathcal{A}$};

% Stati
\node[state]            (q0)  at (0,0)    {$q_0$};
\node[state, accepting] (qi)  at (3,0.8)  {$q_i$};
\node[state, accepting] (qj)  at (3,-0.8) {$q_j$};

% Freccia d'ingresso
\draw[->] ([shift={(-0.7,0.7)}]q0.center) -- (q0);

%%%%%%%%%%%%%%%%%%%% Freccia di trasformazione %%%%%%%%%%%%%%%%%%%%
\draw[->, line width=1pt] (4.3,0) -- node[above] {\textit{Setup}} (6.1,0);

%%%%%%%%%%%%%%%%%%%% Automa trasformato %%%%%%%%%%%%%%%%%%%%
% Rettangolo di A
\draw[-, line width=2.5pt, rounded corners=3pt] (8.6,-1.4) rectangle (11.6,1.4);
\node at (10.1,0) {$\mathcal{A}$};

% Stati
\node[state]            (q0p) at (7,0)      {$q_0'$};
\node[state]            (q0b) at (8.6,0)    {$q_0$};
\node[state]            (qib) at (11.6,0.8) {$q_i$};
\node[state]            (qjb) at (11.6,-0.8){$q_j$};
\node[state, accepting] (qf)  at (13.8,0)   {$q_f$};

% Freccia d'ingresso e epsilon-mossa iniziale
\draw[->] ([shift={(-0.7,0.7)}]q0p.center) -- (q0p);
\draw[->] (q0p) -- node {$\mathbb 1$} (q0b);

% epsilon-mosse verso il nuovo stato finale
\draw[->] (qib) -- node[above, sloped] {$\mathbb 1$} (qf);
\draw[->] (qjb) -- node[below, sloped] {$\mathbb 1$} (qf);
\end{tikzpicture}
\caption{Trasformazione dell'automa $\mathcal{A}$ dopo Setup.}
\label{fig:setup}
\end{figure}

    \item Merge restituisce l'automa in cui due transizioni parallele sono fuse in un'unica transizione: in pratica, poiché l'automa generalizzato è non deterministico è possibile avere due espressioni $e_1$ ed $e_2$ tali che $T(q_i, e_1) = s_1$ e $T(q_i, e_2) = s_2$ e $q_j \in s_1 \cap s_2$. Sono queste due transizioni, dette \textit{parallele} per la coppia di stati $(q_i, q_j)$, che possono essere fuse in un'unica transizione in corrispondenza della somma delle due espressioni regolari.
    
    Definiamo l'automa $(T', q_0, F)$ in cui le transizioni da $q_i$ a $q_j$ sono sostituite con un'unica transizione $T' := \text{Merge}(T, q_i, q_j)$:
    \[
    T'(q, e) = 
    \begin{cases} 
    T(q, e) & \text{se } q \neq q_i \text{ e } T(q, e) \neq q_j \\ 
    \displaystyle \sum_{e \in E_{i,j}} e & \text{se } E_{i,j} := \{e \mid q = q_i \text{ e } T(q_i, e) = q_j\} 
    \end{cases}
    \]

\begin{figure}[H]
\centering
\begin{tikzpicture}[
    >=stealth', shorten >=1pt, auto,
    every state/.style={circle, draw, fill=white, minimum size=0.8cm, inner sep=1pt}
]
%%%%%%%%%%%%%%%%%%%% Automa di partenza %%%%%%%%%%%%%%%%%%%%
% Stati
\node[state] (qi)  at (0,0) {$q_i$};
\node[state] (qj)  at (3,0) {$q_j$};

% Transizioni parallele
\draw[->] (qi) to[bend left=25]  node {$e_1$} (qj);
\draw[->] (qi) to[bend right=25] node[below] {$e_2$} (qj);

%%%%%%%%%%%%%%%%%%%% Freccia di trasformazione %%%%%%%%%%%%%%%%%%%%
\draw[->, line width=1pt] (4.3,0) -- node[above] {\textit{Merge}} (6.1,0);

%%%%%%%%%%%%%%%%%%%% Automa trasformato %%%%%%%%%%%%%%%%%%%%
% Stati
\node[state] (qib) at (7.4,0)  {$q_i$};
\node[state] (qjb) at (10.9,0) {$q_j$};

% Transizione unica
\draw[->] (qib) -- node {$e_1 + e_2$} (qjb);
\end{tikzpicture}
\caption{Fusione di due transizioni parallele da $q_i$ a $q_j$ in un'unica transizione etichettata dalla somma delle espressioni regolari.}
\label{fig:merge}
\end{figure}
\begin{figure}[H]
\centering
\begin{tikzpicture}[
    >=stealth', shorten >=1pt, auto,
    every state/.style={circle, draw, fill=white, minimum size=0.8cm, inner sep=1pt}
]
%%%%%%%%%%%%%%%%%%%% Automa di partenza %%%%%%%%%%%%%%%%%%%%
% Stato
\node[state] (q) at (0,0) {$q$};

% Loop paralleli (archi paralleli da q a q)
\draw[->] (q) edge[loop above] node {$e_1$} (q);
\draw[->] (q) edge[loop below] node {$e_2$} (q);

%%%%%%%%%%%%%%%%%%%% Freccia di trasformazione %%%%%%%%%%%%%%%%%%%%
\draw[->, line width=1pt] (2.0,0) -- node[above] {\textit{Merge}} (3.8,0);

%%%%%%%%%%%%%%%%%%%% Automa trasformato %%%%%%%%%%%%%%%%%%%%
% Stato
\node[state] (qb) at (5.8,0) {$q$};

% Loop unico
\draw[->] (qb) edge[loop above] node {$e_1 + e_2$} (qb);
\end{tikzpicture}
\caption{Caso particolare del Merge con $q_i = q_j = q$: due cappi su $q$ sono transizioni parallele da $q$ a $q$, e vengono fusi in un unico cappio etichettato da $e_1 + e_2$.}
\label{fig:merge-loop}
\end{figure}

    \item Erase Consideriamo un automa generalizzato $(T, q_0, F)$ senza transizioni parallele e introduciamo la seguente trasformazione che consiste nell'eliminare uno stato $\bar{q} \notin F \cup \{q_0\}$ preservando l'insieme che viene deciso dall'automa: il nuovo automa $(T', q_0, F)$ sarà dunque definito su $Q \setminus \{\bar{q}\}$ nel modo seguente: 
    
    se $T(q_i, e_i) = \bar{q}$, $T(\bar{q}, e_j) = q_j$ e $T(\bar{q}, e_k) = \bar{q}$ allora
    \[
    T'(q_i, e) = q_j \text{ dove } e = e_i e_k^* e_j
    \]
    
    per tutte le altre transizioni che non coinvolgono $\bar{q}$ non si interviene e si definisce
    \[
    T'(q, e) = T(q, e).
    \]
\begin{figure}[H]
\centering
\begin{tikzpicture}[
    >=stealth', shorten >=1pt, auto,
    every state/.style={circle, draw, fill=white, minimum size=0.8cm, inner sep=1pt}
]
%%%%%%%%%%%%%%%%%%%% Automa di partenza %%%%%%%%%%%%%%%%%%%%
% Stati
\node[state] (qi) at (0,0) {$q_i$};
\node[state] (qb) at (2.5,0) {$\bar q$};
\node[state] (qj) at (5,0) {$q_j$};

% Transizioni
\draw[->] (qi) -- node {$e_i$} (qb);
\draw[->] (qb) -- node {$e_j$} (qj);
\draw[->] (qb) edge[loop above] node {$e_k$} (qb);

%%%%%%%%%%%%%%%%%%%% Freccia di trasformazione %%%%%%%%%%%%%%%%%%%%
\draw[->, line width=1pt] (6.3,0) -- node[above] {\textit{Erase}} (8.1,0);

%%%%%%%%%%%%%%%%%%%% Automa trasformato %%%%%%%%%%%%%%%%%%%%
% Stati
\node[state] (qib) at (9.4,0)  {$q_i$};
\node[state] (qjb) at (13.4,0) {$q_j$};

% Transizione unica
\draw[->] (qib) -- node {$e_i\, e_k^{*}\, e_j$} (qjb);
\end{tikzpicture}
\caption{Eliminazione dello stato $\bar q$: il cammino $q_i \to \bar q \to q_j$ con cappio su $\bar q$ diventa un'unica transizione etichettata da $e_i e_k^* e_j$.}
\label{fig:erase}
\end{figure}
\end{itemize}
A questo punto
possiamo mostrare il seguente risultato.
\begin{prop}
    Dato $X$ un isnieme regolare, 
    esiste un'espressione regolare 
    $e_X$ tale che
    $\llbracket e_X \rrbracket = X$.
    \begin{proof}
        Consideriamo $A$ un DFA che decide $X$.
        Consideriamo l'automa generalizzato $A'$ definito partendo da $A$
        e considerando come etichette gli elementi formali di 
        del linguaggio $E_A$ invece che le lettere di $A$.
        A questo punto abbiamo un automa generalizzato e possiamo applicare l'operazione "Setup"
        e poi le operazioni "Merge" e "Erase"
        fino a quando non otteniamo un automa con solo due stati, $q_0^*$ e $q_F^*$
        (quelli introdotti con l'operazione "Setup") e un solo arco tra i due stati etichettato 
        con l'espressione $e$.
        Usando che le operazioni sugli automi generalizzati definite prima non modificano l'insieme deciso,
        otteniamo che $e_X=e$.
    \end{proof}
\end{prop}